% Аннотация к ВКР на русском языке, 4–5 страниц без титула — Положение о ГИА Университета Иннополис (2026),
% п. 6.10; титул — Приложение 11. Сборка: make annotation (→ thesis/build/annotation-ru.pdf).
% Текст сжат из глав диссертации (chapters/ru/); правка смысла — сначала в главах, потом здесь.
% Ссылок — только 10 ключевых, чтобы аннотация со списком источников уложилась в 4–5 страниц.
% Результаты пилота — \pending, пока у RQ-1…RQ-4 нет чисел с паспортом (docs/thesis/figures.md).
\documentclass[a4paper,14pt,oneside]{extreport}

\newcommand{\mainlang}{russian}
\newcommand{\docstyle}{thesis}

% Преамбула: шрифты (латиница, кириллица, арабский с огласовками), поля, библиография.
\usepackage{etoolbox}
\usepackage{fontspec}
\usepackage{polyglossia}
\setdefaultlanguage{\mainlang}
\ifdefstring{\mainlang}{english}{\setotherlanguages{russian,arabic}}{\setotherlanguages{english,arabic}}

% Times New Roman есть в macOS; на другой машине — Liberation Serif или Tempora.
\setmainfont{Times New Roman}
\newfontfamily\cyrillicfont{Times New Roman}
% Арабский — тот же шрифт, что рисует карточки бота (ADR-0024): огласовки и шадда не теряются.
\newfontfamily\arabicfont[Script=Arabic,Path=../src/muallim/presentation/telegram/assets/fonts/,
  Extension=.ttf,UprightFont=*-Regular,BoldFont=*-Bold,Scale=1.25]{ScheherazadeNew}
% \ar{…} — арабская вставка в строке; источник строки — в тексте рядом ([SRC-NNN]).
\newcommand{\ar}[1]{\textarabic{#1}}

\usepackage[left=3cm,right=1.5cm,top=2cm,bottom=2cm]{geometry} % [утвердить] по требованиям вуза
\usepackage{setspace}
\onehalfspacing
\usepackage{indentfirst}

\usepackage{graphicx}
\graphicspath{{../docs/architecture/diagrams/}{../docs/design/screens/}{figures/}}
\usepackage{booktabs,tabularx,longtable}
\usepackage{caption}
\usepackage{csquotes}
\usepackage[hidelinks]{hyperref}

% Библиография — тот же sources.bib, что у всего проекта (ключи SRC-NNN).
\ifdefstring{\mainlang}{english}
  {\usepackage[backend=biber,style=ieee,sorting=none]{biblatex}}
  {\usepackage[backend=biber,style=gost-numeric,sorting=none]{biblatex}}
\addbibresource{../docs/research/sources.bib}

% ID рисунков и таблиц из docs/thesis/figures.md: \label{fig:R-01}, \label{tab:T-01}.

% Данные для титульных листов всех документов. [утвердить] — с руководителем и учебным офисом.
\newcommand{\thesisProgramCode}{09.04.01}                     % [утвердить] код направления
\newcommand{\thesisProgramRu}{Информатика и вычислительная техника} % [утвердить]
\newcommand{\thesisProgramEn}{Computer Science}                % [утвердить]
\newcommand{\thesisSpecialtyRu}{Программная инженерия}
\newcommand{\thesisSpecialtyEn}{Software Engineering}
% Тема — предложение автора (2026-10-18; с ADR-0051 — с адаптацией модели), утверждает руководитель [утвердить].
\newcommand{\thesisTopicRu}{ИИ-тренажёр чтения по букварю «Муаллим ас-Сани»: адаптация открытой модели распознавания фонем к речи русскоязычных начинающих и оценка обратной связи}
\newcommand{\thesisTopicEn}{An AI Reading Tutor for the \emph{Muallim al-Sani} Primer: Adapting an Open Phoneme Recognition Model to Russian-Speaking Beginners and Evaluating Pronunciation Feedback}
\newcommand{\thesisShortRu}{ИИ-тренажёр «Муаллим ас-Сани» и адаптация модели}
\newcommand{\thesisShortEn}{AI Tutor for \emph{Muallim al-Sani} and Model Adaptation}
\newcommand{\thesisAuthorRu}{Исмаилов Абдирахим Ерланович}
\newcommand{\thesisAuthorEn}{Abdirakhim Ismailov}               % латиница — [утвердить] как в паспорте
\newcommand{\thesisAuthorEnFull}{Ismailov Abdirakhim Erlanovich}
\newcommand{\thesisSupervisorRu}{Аскарбекулы Нурсултан}
\newcommand{\thesisSupervisorEn}{Nursultan Askarbekuly}
\newcommand{\thesisUniversityRu}{Университет Иннополис}
\newcommand{\thesisUniversityEn}{Innopolis University}
\newcommand{\thesisCity}{Иннополис}
\newcommand{\thesisCityEn}{Innopolis}
\newcommand{\thesisYear}{2027}

\hypersetup{pdftitle={Аннотация: \thesisShortRu},pdfauthor={\thesisAuthorRu}}

% Разделы аннотации — без номеров, в оглавление не идут.
\newcommand{\annsection}[1]{\section*{#1}}

\begin{document}

% ---------------------------------------------------------------- титул (Приложение 11), в объём не входит
\begin{titlepage}
\thispagestyle{empty}
\newgeometry{left=2.5cm,right=2cm,top=2cm,bottom=2cm}
\setstretch{1.0}
\newcommand{\tpbox}[2]{\fcolorbox{black}{softgray}{\parbox[c][#1][c]{\dimexpr\linewidth-2\fboxsep-2\fboxrule}{\centering\bfseries #2}}}
\begin{center}
\textbf{Автономная некоммерческая организация высшего образования\\«Университет Иннополис»}

\vspace{1.2cm}
\textbf{АННОТАЦИЯ\\[4pt]НА ВЫПУСКНУЮ КВАЛИФИКАЦИОННУЮ РАБОТУ\\[4pt](МАГИСТЕРСКУЮ ДИССЕРТАЦИЮ)\\[4pt]
ПО НАПРАВЛЕНИЮ ПОДГОТОВКИ\\[4pt]\thesisProgramCode{} \MakeUppercase{\thesisProgramRu}}

\vspace{1cm}
\textbf{НАПРАВЛЕННОСТЬ (ПРОФИЛЬ) ОБРАЗОВАТЕЛЬНОЙ ПРОГРАММЫ\\[4pt]«\thesisSpecialtyRu»}
\end{center}

\vspace{0.8cm}
\noindent\begin{tabularx}{\linewidth}{@{}P{2.6cm}X@{}}
\textbf{Тема} & \tpbox{2.4cm}{\thesisTopicRu}\\
\end{tabularx}

\vspace{0.8cm}
\noindent\begin{tabularx}{\linewidth}{@{}P{2.6cm}XP{3.2cm}@{}}
Выполнил & \tpbox{1.7cm}{\thesisAuthorRu} &
  \fcolorbox{black}{softgray}{\parbox[b][1.7cm][b]{2.9cm}{\centering\tiny подпись}}\\
\end{tabularx}

\vfill
\begin{center}\thesisCity, \thesisYear\end{center}
\restoregeometry
\end{titlepage}
\setcounter{page}{2}

% ---------------------------------------------------------------- введение, цель и задачи
\annsection{Введение, цель и задачи}

Взрослые, которые учатся читать Коран, начинают с букваря и учителя. В России и Центральной Азии таким
букварём больше века остаётся «Муаллим ас-Сани» Ахмад-Хади Максуди: он перенёс звуковой метод джадидов на
чтение Корана --- буква изучается как звук, слоги строятся из неё малыми шагами \cite{SRC-015}.
Книгу и сейчас издаёт Духовное управление мусульман Татарстана и официально распространяют в Казахстане \cite{SRC-040}. Урок идёт по циклу \emph{талакки}: учитель читает, ученик повторяет, учитель
исправляет каждую ошибку \cite{SRC-021}. Цикл работает, но привязывает каждого ученика ко времени
учителя: между уроками взрослый читает один, и никто не скажет ему, что \ar{ح} превратилась в русское [х].

Открытая модель Quran Muaalem распознаёт фонемы коранического арабского с частотой ошибок около половины
процента, но обучена только на чтецах-экспертах; на записях с реальными ошибками её ошибка в десять раз выше \cite{SRC-002}. Публичные наборы для оценки чтения Корана не содержат учеников, для которых арабский не
родной \cite{SRC-004}. Кроме того, модель оценивает одну попытку, а учителя интересует, что
ученик \emph{уже усвоил}. Обе проблемы упираются в отсутствие набора данных: без него нельзя ни измерить, как
модель ошибается на начинающих, ни исправить её, ни проверить оценку усвоения.

\textbf{Объект исследования} --- набор данных букварного чтения взрослых русскоязычных начинающих и ИИ-модель,
которая распознаёт в нём ошибки чтения и оценивает усвоение навыков. \textbf{Цель} --- разработать такой
набор данных и модель и измерить, насколько модель приближается к преподавателю и что она показывает о том,
как учатся люди.

\textbf{Задачи:}
\begin{enumerate}[label=З\arabic*.]
  \item перевести букварь в цифровую программу в порядке книги с привязкой к странице и строке издания и
    построить карту навыков;
  \item создать тренажёр --- инструмент сбора данных, который распознаёт фонемы, объясняет ошибки по-русски и
    сохраняет полный след решения;
  \item собрать набор данных с согласия участников и разметкой двух преподавателей по фонемам и типу ошибки;
  \item адаптировать модель распознавания на поправках учителей и минимальных парах книги;
  \item построить модель усвоения навыков книги с учётом измеренных ошибок распознавателя;
  \item описать продвижение учеников по навыкам книги и ответить на исследовательские вопросы.
\end{enumerate}

Исследовательские вопросы: насколько надёжно учителя и модель различают ошибки чтения и произношения (RQ-1);
насколько адаптация улучшает распознавание на новых учениках и как это зависит от объёма данных (RQ-2);
можно ли по распознаванию оценить усвоение каждого навыка и улучшает ли адаптация эту оценку (RQ-3); как
русскоязычные взрослые продвигаются по навыкам книги (RQ-4).

% ---------------------------------------------------------------- обзор литературы
\annsection{Обзор литературы}

Обзор охватывает пять областей. \emph{Букварь и его метод}: звуковой пошаговый метод с немедленным
исправлением хорошо сочетается с автоматической обратной связью, а уроки книги дают готовую карту навыков. \emph{Компьютерное обучение произношению}: метаанализ показывает средний
положительный эффект автоматической обратной связи, наибольший --- когда она касается отдельных звуков \cite{SRC-025}; при этом распознаватели дают ложные срабатывания на речи с акцентом.
\emph{Автоматическая оценка чтения Корана}: Quran Muaalem распознаёт фонемы вместе с артикуляционными
признаками \cite{SRC-002}, но бенчмарки содержат беглое чтение аятов, а не букварное чтение начинающих \cite{SRC-004}. \emph{Адаптация речевых моделей}: дообучение низкого ранга \cite{SRC-111},
калибровка и отказ от ответа позволяют адаптировать модель на малых данных, но требуют
данных новой аудитории. \emph{Репетиторы чтения и отслеживание знаний}: системы, которые слушают чтение вслух,
работают для детей \cite{SRC-125}, а отслеживание знаний оценивает усвоение по
последовательности попыток \cite{SRC-124}, в том числе по шумному выходу распознавателя \cite{SRC-125}.

Пробел состоит из четырёх пунктов: нет набора чтения взрослых русскоязычных начинающих, где ошибки чтения
отделены от ошибок произношения; неизвестно, можно ли на малом наборе этой аудитории адаптировать модель так,
чтобы она ближе подошла к учителю; модели усвоения по речи не учитывают ошибки распознавателя, измеренные
против учителей; не измерено, как взрослые русскоязычные осваивают навыки арабского чтения.

% ---------------------------------------------------------------- основное содержание
\annsection{Краткое содержание работы}

\textbf{Методология.} Работа выполнена в парадигме design science research \cite{SRC-100}: знание
получают, создавая артефакты и строго их оценивая. Артефакт A --- набор данных с инструментом сбора, B ---
модель распознавания, которая учится у учителей, C --- модель усвоения навыков книги. Требования записаны как
сценарии атрибутов качества, метрики выведены методом «цель --- вопрос --- метрика». Гипотезы, разбиение данных и план анализа предрегистрированы и зафиксированы отпечатком до
сбора данных и до первого обучения. Модели сравниваются на замороженном тесте из учеников, не участвовавших в
обучении, со слепой разметкой.

\textbf{Реализация.} Инструмент сбора --- Telegram-тренажёр, который ведёт ученика по всей книге в её
порядке: 125 уроков и 800 карточек, каждая прослеживается до страницы и строки издания 2010~г. Ученик
присылает голосовое, Quran Muaalem распознаёт фонемы, ученик получает объяснение на русском, а
попытка с полным следом решения сохраняется в PostgreSQL; аудио хранится только с согласия и под псевдонимом.
Два преподавателя размечают попытки по фонемам и различают ошибки чтения (прочитан не тот знак) и ошибки
произношения (знак верный, звук неверный). На собранных данных строятся калиброванный верификатор с исходом
«не уверен» и гарантией на долю ложных отклонений и дообучение LoRA с пополнением по минимальным парам книги, а
поверх них --- байесовское отслеживание знаний по навыкам книги с учётом измеренных ошибок распознавателя.
Система создана автором вместе с ИИ-агентом под механическими проверками, которые защищают и код, и
эксперименты.

\textbf{Оценка.} Предварительная проверка показала, что конвейер отличает верное чтение эксперта от
неверного (99--100 против 40--67 баллов) и отвечает быстрее 10 секунд. \pending{результаты пилота, адаптации
и оценки усвоения по RQ-1--RQ-4}.

\textbf{Новизна.} Набор данных --- первое собрание чтения взрослых русскоязычных начинающих, привязанное к
навыкам букваря, с разделением ошибок чтения и произношения; модель усвоения учитывает ошибки распознавателя,
измеренные против учителей; кривые обучения впервые описывают, как взрослые русскоязычные продвигаются по
навыкам «Муаллим ас-Сани».

% ---------------------------------------------------------------- заключение
\annsection{Заключение}

Работа даёт четыре результата: набор данных букварного чтения взрослых русскоязычных начинающих с разметкой
учителей и карточкой набора (C1); адаптацию открытой модели к новой аудитории на малых данных с кривой обучения
и вкладом каждого источника данных (C2); оценку усвоения по речи с учётом ошибок распознавателя (C3); кривые
обучения взрослых по навыкам букваря (C4). Инженерный результат --- воспроизводимый конвейер «сбор $\to$
разметка $\to$ обучение $\to$ оценка» с заменяемыми моделью, базой данных и мессенджером.
\pending{ответы на RQ-1--RQ-4}.

Магистерская работа задумана как первый этап докторского исследования: выбор следующей карточки по усвоению, активное обучение на разметке учителей, контролируемый эксперимент с обратной связью и перенос
на другие родные языки и правила таджвида. Тренажёр задуман как дополнение к учителю, а не замена ему.

% Раздел, а не глава: без разрыва страницы, чтобы уложиться в 4–5 страниц.
\annsection{Список использованных источников}
\printbibliography[heading=none]
\end{document}
